\section{Pre-registration, addenda and validation gates}
\label{supp:prereg}

This section carries the documentary basis of the manuscript's
pre-registration claim. It reproduces, in the order in which they were
written, (i) what was registered on 2026-08-25 before any model was fit to
real Human Mortality Database data, (ii) the four registered addenda, each
hash-stamped in its own commit, with the SHA-256 digest of every document and
a clause-by-clause account of what each addendum fixed, and (iii) the three
validation gates that had to pass before any real-data number could be read,
including the synthetic-truth design, the \pkg{StMoMo} oracle-parity result
and the life-table parity check.

It supports two statements the manuscript makes and does not have room to
document. The first is the integrity claim of Section~4: none of the four
addenda changes a population list, a split year, a horizon, a test-window
observation or the metric list. A reader can check that claim only against
the addenda themselves, which are therefore given here in substance rather
than summarised. The second is the harness-validity claim: every coverage
number in the paper is a statement about a forecasting method only if the
evaluation code measures coverage correctly, and Section~\ref{supp:gates}
below is the evidence that it does. The deviations from the registered design
are \emph{not} here: they stay in the manuscript, in Section~4, because they
are findings rather than housekeeping.

\subsection{What was registered, and when}
\label{supp:prereg-what}

The registration is the file \texttt{PREREGISTRATION.md}, committed
2026-08-25, before any model had been fit to real HMD data. It pins the data
vintage (HMD bulk files, last modified 2026-06-15, Methods Protocol~v6,
checksummed in \texttt{data/MANIFEST.sha256}), the global seed 20260825 with
per-ensemble-member seeds recorded in the results metadata, and the following
analytic choices.

\paragraph{Question.} Do the prediction intervals of stochastic,
machine-learning and neural mortality forecasting models attain their nominal
coverage when the evaluation window crosses a structural break (COVID-19),
and does the answer depend on the model family, the uncertainty mechanism, or
both?

\paragraph{The five hypotheses, as registered.} These are the sentences fixed
before estimation; Addendum~3~\S8 later recorded how the prior-art map
changes the reading of each, and the manuscript reports the resulting
deviations in the main text.

\begin{itemize}
\item \Hone: model rankings by RMSE and by CRPS disagree (rank correlation
  below one; at least one top-three inversion) in the shift regime.
\item \Htwo: nominal 95\% intervals under-cover in the shift regime for every
  model family; the shortfall exceeds the stable-regime shortfall.
\item \Hthree: joint path coverage over $h=1,\dots,5$ is materially below
  marginal coverage at the same nominal level for every family (reported gap;
  test: paired difference greater than zero).
\item \Hfour: coverage failure is non-uniform in age; shift-regime coverage at
  ages 65--99 is lower than at ages 25--64 for every family.
\item \Hfive: interval miscalibration propagates to derived quantities; the
  empirical coverage of 95\% intervals for $e_{65}$ and the annuity factor
  $\annuity$ at 2\% departs from nominal by more than the corresponding
  stable-regime departure.
\end{itemize}

No directional claim was registered about neural-versus-classical ordering;
the crossed design estimates it.

\paragraph{Data.} HMD $1\times1$ deaths and central exposures, ages 0--99
(the registration already records that cells with $E=0$ are all at ages 90 and
above in this vintage, and excludes the $110{+}$ open group and the
model-smoothed old ages), sexes modelled separately, Poisson likelihood with
log-exposure offset throughout; rates are never modelled as Gaussian on the
raw scale.

\paragraph{Regimes and populations.} \emph{Shift} (primary): train to 2019,
test 2020--2024, $h=1,\dots,5$, on the twenty populations whose final annual
data reach 2024 in this vintage (BEL, CHE, CHL, DNK, EST, FIN, HKG, HRV, ISL,
JPN, KOR, LTU, LUX, LVA, NOR, PRT, SVK, SWE, TWN, USA). \emph{Stable}
(control): expanding origin, train to $T$ for $T=1990,1992,\dots,2014$, test
$T+1,\dots,T+5$ with all test years at or before 2019, the same twenty
populations subject to each population's own start year. \emph{Placebo
break}: train to 1913, test 1914--1922, $h=1,\dots,9$, on the eleven
populations with continuous 1908--1922 coverage and a start year no later
than 1900 (CHE, DNK, FIN, FRATNP, GBRTENW, GBR\_SCO, ISL, ITA, NLD, NOR,
SWE), Belgium excluded for its missing occupation-era rows and FRACNP and
GBRCENW excluded as overlapping variants of FRATNP and GBRTENW.

\paragraph{Factors.} Ten model families (Lee--Carter by SVD with a random
walk with drift; Poisson Lee--Carter; CBD~M5; APC/Renshaw--Haberman~M2-A;
sparse VAR on log rates; multi-output Gaussian process; neural-LC with
Richman--Wuethrich embeddings; shallow CNN-LC; LSTM on $\kappa_t$;
distributional Poisson/NB-head network) crossed with seven uncertainty
mechanisms (model-native; semiparametric Poisson bootstrap; deep ensemble
with $M=10$; Monte Carlo dropout; split conformal; EnbPI; copula/joint
conformal). The registration states in advance that not every cell is
admissible, enumerates the admissible sub-grid in \texttt{docs/GRID.md} at 50
primary cells, and requires every claim to be stated as a contrast within an
admissible sub-grid, never as a full-factorial main effect. Conformal
calibration sets are drawn only from data available at the training cutoff.

\paragraph{Metrics.} All computed from predictive samples through one code
path: RMSE and MAE on $\log \hat m_x$ and the $e_0$ error; CRPS by the sample
estimator on $\log m_x$; the Poisson log score on observed deaths rounded to
the nearest integer, with CRPS on death counts as the registered sensitivity;
Winkler/interval score and empirical coverage at 50/80/95\%; randomised PIT
with a uniformity test; calibration-by-age curves; joint path coverage over
$h$ per population--sex--age trajectory; the Murphy
calibration--resolution decomposition; interval coverage and error for
$e_{65}$, $e_0$ and $\annuity$ at 2\% interest from each predictive sample's
period life table; and, for inference, Diebold--Mariano per population--sex on
CRPS with a wild cluster bootstrap over populations and a model confidence
set at 90\%.

\paragraph{Splits discipline and planned robustness.} Expanding-origin time
splits only; hyperparameters tuned on inner time splits (the final five
pre-cutoff years of each training window); the test years touched exactly
once, by the final scoring run. The registered secondary runs are STMF weekly
data for break timing, sensitivity to 2024-provisional revisions (re-run the
shift regime dropping 2024) and sensitivity to an age cap of 90 rather
than~99.

\subsection{The hash ledger}
\label{supp:prereg-hashes}

Each document below was committed before the work it governs and stamped with
the SHA-256 digest of its own bytes, so any later edit is detectable. In the
order registered:

\begin{itemize}
\item \texttt{PREREGISTRATION.md}, 2026-08-25, before any real-data fit:\\
  \texttt{\small bbbe3860a5446d063e186e56555afa893b07c36cf15a990d07567471affe50be}
\item \texttt{PREREGISTRATION-ADDENDUM-1.md}, 2026-08-26, commit
  \texttt{5f2e3cb}, still before any real-data fit:\\
  \texttt{\small 8b43083afae9f0eaafce5e89381349d74cbbafb94022b56d1f2f39032a308889}
\item \texttt{PREREGISTRATION-ADDENDUM-2.md}, 2026-08-26, commit
  \texttt{577ce52}, still before any real-data fit:\\
  \texttt{\small 1618e9fdcc1d6cbc760ce40ea496c1285fe44b8fc7520e8687e2101343ebdbc1}
\item \texttt{PREREGISTRATION-ADDENDUM-3.md}, 2026-08-26, commit
  \texttt{01f154a}, still before any real-data fit:\\
  \texttt{\small cca2881a3a7ae6a17bbec7a4c6440192a6f9aa1cdd5766c13ef0aaec8ddeb16f}
\item \texttt{PREREGISTRATION-ADDENDUM-4.md}, 2026-08-29, during the
  production sweeps and before any GP cell on a panel longer than sixty years
  had produced a result:\\
  \texttt{\small 586e01780cceeea32878d1a493662c2883d017396459dc3442a6053f101a564b}
\end{itemize}

Addendum~3 additionally cross-references the prior-art map
\texttt{docs/PRIOR-ART.md} (SHA-256
\texttt{\small 3f36fb62615deedb1dea9505e313581ace8ed3be3eb4962c63cdfd58164bdd12}),
committed immediately before it, whose verdicts drive the hypothesis-status
clauses of \S8 of that addendum.

None of the four addenda changes a population list, a split year, a horizon, a
test-window observation or the metric list. Addenda~1, 2 and~3 were written
before any model had been fit to real data at all; Addendum~4 was written
during the sweeps, and its scope, cause and disclosure are set out in
Section~\ref{supp:add4} and, as a deviation, in the manuscript.

\subsection{Addendum 1: data-prerequisite strata and sensitivities}
\label{supp:add1}

Addendum~1 arises from the HMD documentation check recorded in
\texttt{docs/DATA-PREREQS.md} and adds pre-specified sensitivity analyses
only.

For the \emph{placebo} regime the primary panel is unchanged. Belgium's
exclusion is confirmed by HMD's own data-series gap warning for 1914--1918,
and FRACNP is never used, not even as a sensitivity, because HMD flags its
1919 civilian denominators as implausible. The eleven placebo populations are
stratified from HMD's country documentation, and the strata are reported
alongside the pooled placebo result: \emph{neutral, register-based} (CHE, DNK,
FIN, ISL, NLD, NOR, SWE), whose only break in the window is the 1918
influenza pandemic and which is therefore the clean analogue for the
1918-versus-2020 comparison; \emph{belligerent, total series} (FRATNP,
GBRTENW, ITA), in which HMD reconstructs military deaths at male ages of
roughly 18--45; and \emph{civilian-only} (GBR\_SCO), in which deaths abroad
are excluded by construction. Three secondary sensitivities are pre-specified:
drop GBR\_SCO, whose 1912--1938 population rests on an ad-hoc spline HMD
itself calls unreliable; report the neutral stratum alone; and report Denmark
with and without 1921--1922, the 1921 South Jutland territorial change falling
inside the test window.

For the \emph{shift} regime the addendum records that HMD carries no per-year
provisional flag and re-issues whole series at each release, pins the vintage
by DOI 10.4054/HMD.Countries.20260615, and names the known denominator risks:
USA 2020--2025 exposures blend preliminary 2020-census inputs, CHL exposures
rest on the 2017 census pending revision, and TWN 2024 deaths are provisional
per STMF, while the register-based populations (the Nordic countries, CHE,
NLD, EST, LVA, LTU) are effectively final. Beyond the already-registered
drop-2024 run it adds three sensitivities: drop USA and CHL entirely; a
register-based versus census-based contrast; and a re-run on the next HMD
vintage if one is released before submission, with any change in headline
coverage reported as a vintage-revision effect.

The addendum also records that HMD's Lexis-triangle split for 1918 (Methods
Protocol~v6, Appendix~A) affects only how raw $1\times1$ death squares are
apportioned into triangles and leaves \texttt{Deaths\_1x1} untouched where the
raw data were already $1\times1$: a data-appendix sentence with no modelling
consequence.

\subsection{Addendum 2: scoring-target clarifications}
\label{supp:add2}

Addendum~2 clarifies how two registered metrics are computed and changes no
hypothesis, population, split, horizon or metric.

\begin{enumerate}
\item \emph{Rate-scale scores evaluate the predictive law of the observed
  rate.} The registration says CRPS, Winkler scores, coverage and PIT are
  computed ``on $\log m_x$''. The observed $m_x=D/E$ carries Poisson sampling
  noise that a model's latent-rate samples do not, so scoring latent samples
  against observed rates would count observation noise as miscalibration,
  most severely at young ages and in small populations. Rule: every
  rate-scale score is computed on Poisson-inclusive predictive samples,
  $D^{\ast}\sim\mathrm{Poisson}(E\,m^{\ast})$ drawn on the model's own rate
  paths and converted to log crude rates---the same construction already
  registered for the count-scale metrics. Native, bootstrap, ensemble and
  conformal mechanisms are treated identically, so the crossed design is
  unaffected.
\item \emph{Zero-death cells: half-count continuity correction, no masking.}
  Log rates, observed and sampled alike, use $\max(D,\tfrac12)/E$, the
  standard demographic continuity correction applied to both sides of every
  score. Cells are not dropped; the number of zero-death cells per
  (population, sex, regime) is reported. Count-scale metrics need no
  correction.
\item \emph{Conformal cells are scored on their intervals, at their
  construction level.} Coverage and Winkler scores for conformal cells come
  from the interval bounds themselves at the construction level of 95\%; the
  50\% and 80\% columns are reported as not applicable for those cells; and
  CRPS, log score and PIT for them remain flagged secondary. This removes a
  scoring bias against conformal arms that taking empirical quantiles of
  uniform-in-interval samples would otherwise create.
\item \emph{PIT uniformity: statistic and a caveated $p$-value.} The
  Kolmogorov--Smirnov statistic is reported with its nominal $p$-value, which
  is descriptive because PIT values across ages and horizons within a
  population are dependent. Formal inference on calibration uses the
  population-clustered procedures already registered.
\end{enumerate}

\subsection{Addendum 3: defect handling, scoring repairs and hypothesis status}
\label{supp:add3}

Addendum~3 was registered before any model had been fit to real data and
before any real-data result existed. Sections~1--7 and 9--11 specify
computation on cases the original text already covers; \S8 records
hypothesis-status clarifications forced by the prior-art map, made while no
result existed. Each clause is cited in the manuscript at the point where it
applies.

\paragraph{\S1 Zero-exposure training cells: $W=\ind\{E>0\}$.} An audit of the
2026-06-15 vintage found 99 cells with $E=0.00$, every one with $D=0.00$, at
ages 95--99 in seven (population, sex) series of the shift training panel
(CHE female 1, CHE male 10, FIN male 27, ISL female 13, ISL male 30, LUX male
6, SWE male 12), verified as structural zeros: HMD's 1~January population
count is 0.00 at both corners of every one of the 99 Lexis squares. Under the
registered Poisson likelihood such a cell contributes an identically zero
log-likelihood, score and Fisher information, so the maximum-likelihood
estimate with the cell included equals the estimate under indicator weights
exactly, and only the numerical evaluation of $0\log 0$ fails. Every family
therefore carries cell weights $W=\ind\{E>0\}$, \pkg{StMoMo}'s \texttt{wxt}
convention: the numerically correct evaluation of the registered objective,
recorded as a clarification and not a deviation. The per-family
specifications, and the measured row losses for the sparse VAR's
common-complete-subsample rule, are reproduced in
Section~\ref{supp:methods-zeros}.

\paragraph{\S2 Training-year contiguity; Belgium.} The training window for
every (population, sex) is the maximal contiguous block of complete years
ending at the origin. The binding case is BEL, whose 1914--1918 deaths and
exposures are missing at every age, so BEL trains on 1919--2019 in the shift
regime. BEL remains excluded from the placebo regime as already registered.

\paragraph{\S3 Zero-exposure cells in a test window.} Scoring masks ages
exactly as the runner's age mask does---an age is scored only if every
horizon's observation is finite---with \texttt{n\_ages\_scored} reported and
\texttt{cov95\_by\_age} null on masked ages. Derived quantities are computed
on the maximal contiguous scored age range starting at age~0, on the
predictive and observed sides symmetrically, with the range reported. The
binding cells are in the placebo test window at ages 98--99 (CHE male~1, FIN
male~3, ISL female~5, ISL male~1).

\paragraph{\S4 Minimum training length.} A (origin, population, sex) triple is
admissible only with at least fifteen complete training years; inadmissible
triples are recorded as such and not scored, and the per-origin effective
cluster count is reported alongside every stable-regime summary (the stable
panel is unbalanced, growing from 27 to 33 buildable population--sex cells
between origins 1990 and 2014). The measured motivation is recorded in the
addendum: at origin~1994 with three training years, CHL males wrote an
LC/native \texttt{coverage\_95} of 0.784 as an apparently valid row.

\paragraph{\S5 Rate-scale truth uses rounded deaths.} The rate-scale observed
truth is $\log(\max(\text{round}(D),\tfrac12)/E)$, so that both sides of every
rate-scale score live on the same lattice. The measured motivation: scoring
fractional Lexis-split deaths against an integer lattice moved the PIT KS
statistic from 0.019 to 0.133 at $\lambda=2$ and cost 2.2 percentage points of
95\% coverage at $\lambda=5$---a false-positive generator aimed squarely at
\Htwo\ and \Hfour. HMD death counts are non-integer for nine of the twenty
shift populations and exactly integer for the other eleven, and the
per-population non-integer share is reported. The observed \emph{life table}
of \Hfive\ keeps raw fractional deaths and a $10^{-10}$ rate floor, because at
an age where nobody died $m=0$ is the correct life-table statement, while the
half-count would shift ISL female 2020 $e_0$ by $-0.370$ years, beyond
gate~3's 0.15-year tolerance. The two conventions in the same cell are
deliberate.

\paragraph{\S6 Conformal cells: scored from their own bounds.} Each wrapper
exposes its interval at its own construction level (the wrapper's $\alpha$,
not a hard-coded 0.95), and coverage and Winkler scores are computed from
those bounds directly, with no uniform-in-interval sampling and no Poisson
composition---the radius is calibrated on observed residuals and already
contains observation noise, and composing again was measured to inflate width
by 150.7\% at $\lambda=10$. Conformal calibration residuals use the same
half-count convention as the scorer, replacing a rate floor under which a
single zero-death calibration cell inflated a SWE female band radius from
1.44 to 12.94 nats; the conformal quantile ignores non-finite residuals.

\paragraph{\S7 Sparse VAR: explosive draws are rejected, not clipped.} A
coefficient draw whose companion-matrix spectral radius is at least one is
rejected and redrawn, capped at one hundred attempts per path; a path that
cannot produce a stable draw makes the cell an error row. Predictive rates and
Poisson means are never clipped to rescue a divergent draw, because clipping
would convert divergence into a bounded interval and silently flatter the
family's coverage. Without the rule, real panels produced predictive $m_x$ up
to $3\times10^{45}$ and a hard failure in the Poisson composer.

\paragraph{\S8 Hypothesis-status clarifications.} Recorded from the prior-art
map while no real-data result existed. \Hone\ is demoted from headline
hypothesis to harness-validity check, the point-versus-proper-score ranking
divergence being already published; the residual claim is its magnitude under
pre-registration, per population and unweighted by exposure. \Htwo\ is
restated two-sided: the tested quantity is
$|\text{empirical coverage}-\text{nominal}|$, hypothesised to increase from
the stable to the shift regime with direction free to differ by (family,
mechanism), the registered one-sided form being already falsified in the
stable regime by published neural intervals. \Hthree's comparator is
re-specified to each model's own simulated-path-implied joint coverage,
computable from the same predictive sample paths, because the registered
comparison with the marginal rate is bounded toward confirmation
($0.95^5=0.774$ under independence) and would be near-unfalsifiable; the
realised-minus-implied gap can be of either sign. \Hfour's registered
direction is contradicted by \citet{dowd2010backtesting}, and the addendum
pre-commits to reporting a reversal as an informative result rather than
re-testing until confirmation. \Hfive\ is restricted to the ex-post coverage
of derived-quantity intervals, the construction of such intervals being prior
art, with \citet{cairns2011mortality}'s attenuation finding as the named
alternative.

\paragraph{\S9 Conformal calibration horizon covers the placebo.} The
calibration horizon must be at least the longest scored horizon in every
regime, so the wrappers' calibration horizon moves from eight to nine and no
placebo horizon reuses a radius calibrated for a shorter one. The shift regime
is unaffected.

\paragraph{\S10 Zero-death-cell reporting.} The count is per (population, sex,
regime), computed on the full age range 0--99 from the observed test window
before any model age mask, with the criterion $D<0.5$---the threshold at which
the half-count correction binds, and the criterion under which HMD's
fractional deaths make a test of $D=0$ under-count by 43\% in the placebo. The
corrected figures for the 2020--2024 window are ISL female 117 of 500 (23.4\%),
ISL male 68 of 500, LUX female 110 of 500 (22.0\%) and LUX male 72 of 500; the
addendum records that a previously cited 52 cells (10.4\%) had been computed
on HMD's \emph{Total} column, which the study never models.

\paragraph{\S11 Contrasts on common cells.} Conformal arms fail on short
panels and native arms fail where fits diverge---different cells, by
mechanism---so any contrast between mechanisms or families is computed on the
intersection of (origin, population, sex) cells in which every compared arm
produced a valid row, with the intersection size reported and the full
uncensored tables carried in the supplement. Error rows are never silently
dropped from a denominator: every table reports its cell count.

\subsection{Addendum 4: the sixty-year training-window cap for the GP}
\label{supp:add4}

Addendum~4 is the one addendum registered during the production sweeps, on
2026-08-29, before any GP cell on a panel longer than sixty years had produced
a result. It applies to the multi-output Gaussian process under every
mechanism in every regime.

\emph{Cause.} The exact multitask GP's kernel memory scales with
$(\text{years}\times\text{ages})^2$: a 269-year panel is 26{,}900 points and a
5.8~GB kernel per fit, and the EnbPI and copula-conformal mechanisms refit the
family ten times per cell. On the compute node the second-pass GP sweep
stalled after nineteen hours with one worker holding 10.7~GB, the machine
swapping, and fourteen of twenty populations untouched; six had completed
(CHL, HKG, TWN, LUX, and KOR and HRV as design-floor errors). The arm as
specified is infeasible on long panels on the available hardware.

\emph{Rule.} The GP family trains on the trailing sixty complete years of each
training panel (\texttt{max\_years = 60}), applied through the runner's
per-family keyword arguments so that every mechanism sees the same family.
Panels shorter than sixty complete years are unaffected.

\emph{Why sixty, and why this is not outcome-driven.} Sixty is the longest
panel the arm completed without a cap (LUX), so the rule changes nothing
already computed and is defined by feasibility rather than by results; the
four populations already completed have panels of 28--60 years and their cells
are identical under the rule, so they are kept. Sixty is also more than twice
the 27-year window used by the family's source paper
\citep{huynh2021multi}, which the pre-registration cites as the reference
specification. Disclosure, recorded in the addendum itself: before any cap was
contemplated, a single diagnostic on one population had shown a trailing
forty-year window scoring \emph{better} than the full 269-year panel
(\texttt{docs/STATUS.md}, 2026-08-27); the cap chosen here is sixty, not
forty, it was not selected by score, and the earlier observation is recorded
so that the reader can judge. No other family, mechanism, regime, metric,
hypothesis or inference rule changes, and the results tables state the window
for the GP family.

\subsection{Validation gates}
\label{supp:gates}

Three gates had to pass before any real-data result could be read. All three
passed, and gates~1 and~2 were re-run after the weighted-fit rewrite of
Addendum~3~\S1. If gate~1 had failed, every coverage number in this paper
would have been measuring a bug rather than a forecasting method.

\paragraph{Gate 1: synthetic truth.} Data are simulated from a known Poisson
Lee--Carter process---40 ages, 60 training years, horizons $1$--$10$,
exposures $10^5$, thirty independent worlds, 1{,}500 predictive samples per
world---and pushed through the same code path the real-data runs use. The
gate is a confidence-interval test in
\texttt{tests/test\_synthetic\_calibration.py}: the correctly specified
Poisson Lee--Carter must attain nominal 95\% coverage within Monte Carlo
tolerance and produce a near-uniform PIT histogram, and the harness must
detect an injected break as under-coverage. It does both.

Each neural and GP family passed its own sub-gate on the same synthetic
worlds: finite and bounded point error over $h=1,\dots,5$; distinct recorded
member seeds for every ensemble; and, for the distributional NB head,
reproduction of an NB-simulated world's coverage through the Poisson
composition of Section~\ref{supp:methods-interface}. One documented finding
came out of the gate and is carried into the interpretation of every SVD-LC
row in the paper: the two-stage SVD Lee--Carter is mildly \emph{over}-dispersed
on Poisson data, because the SVD $\kappa_t$ carries observation noise into the
random-walk innovation variance, so its intervals are slightly wide and its
PIT centre-heavy. This is a property of the method, not of the harness.

\paragraph{Gate 2: oracle parity against R \pkg{StMoMo}.} The Python Poisson
Lee--Carter must match R \pkg{StMoMo} (version 0.4.1 under R~4.6.1) on a
shared subset. On Swedish males, 1950--2000, ages 0--99, the two fits have
\emph{identical} Poisson log-likelihood and maximum relative differences of
$3.1\times10^{-15}$ in $\alpha_x$, $2.4\times10^{-14}$ in $\beta_x$ and
$1.5\times10^{-13}$ in $\kappa_t$ when the Newton iteration is run to its cap.
The qualifier matters and is stated rather than suppressed: the agreement
quoted is the fixed point of the iteration, not an early-stopped fit. The
\pkg{StMoMo} parameter dump is tracked in the repository, so the check can be
re-run without an R installation.

\paragraph{Gate 3: life-table parity.} Our period life table computed on HMD's
published $m_x$ must reproduce HMD's published $e_0$ and $e_{65}$ within
0.15~years at spot checks. It does. The tolerance is the size of the
documented $a_x$ simplifications between our routine and HMD's, which are set
out in Section~\ref{supp:methods-lifetable}.
